\paragrafos{Διαίρεση ρητών: πόσες φορές χωράει;}{Α.3.5, «Διαίρεση ρητών»}

\begin{anagke}
Ο παππούς Γιάννης έχει $450$ λίτρα ελαιόλαδο και θέλει να το βάλει σε δοχεία που χωράνε από $1\frac{1}{2}$ λίτρα. «Πόσα δοχεία θα χρειαστώ;» ρωτάει. Ο Ηλίας γράφει $450:1\frac{1}{2}$ και κολλάει: «Πώς διαιρώ με μεικτό;». Ο Στέλιος προσπαθεί να εφαρμόσει έναν κανόνα που μόλις έμαθε: «Διαιρώ τους αριθμητές μεταξύ τους και τους παρονομαστές μεταξύ τους». Τον σταματά η Ρήνη: «Σε ένα δοχείο χωράνε λίγο παραπάνω από ένα λίτρο. Τα δοχεία δεν μπορεί να είναι περισσότερα από $450$». Η Ζωή ρωτάει «πόσες φορές χωράει;». Μετράνε: κάθε $3$ λίτρα γεμίζουν $2$ δοχεία, άρα τα $450$ λίτρα γεμίζουν $300$ δοχεία.

Η Πολυξένη παρατηρεί ότι το $300$ προκύπτει και ως $450\cdot\frac{2}{3}$. «Το να διαιρέσεις με το $\frac{3}{2}$ είναι το ίδιο με το να πολλαπλασιάσεις με το $\frac{2}{3}$, τον αντίστροφό του;» Η Αγγελική προσθέτει ένα δεύτερο πρόβλημα: «Στο ταμείο της εκδρομής έλειπαν $24$ ευρώ. Αν χωριστούν σε $6$ παιδιά, πόσο έλειπε σε καθένα;» Και ο Νίκος: «Και αν ο διαιρέτης είναι δεκαδικός, π.χ. $3{,}6:0{,}45$;».

Σε αυτή την παράγραφο θα μάθουμε να διαιρούμε κάθε είδους ρητούς: ακεραίους, κλάσματα, δεκαδικούς και θετικούς και αρνητικούς.
\end{anagke}

\begin{lexi}{διαίρεση, πηλίκο, λόγος}{διαίρεση, πηλίκο, λόγος}
Η \lexiword{διαίρεση} είναι το «δι-αιρώ», «χωρίζω σε μέρη». Ο \lexiword{διαιρετέος} είναι αυτός που μοιράζεται, ο \lexiword{διαιρέτης} αυτός που μοιράζει και το \lexiword{πηλίκο} (από το «πηλίκος», «πόσος;») το αποτέλεσμα. Το πηλίκο $\alpha:\beta$ ή $\frac{\alpha}{\beta}$ λέγεται και \lexiword{λόγος} του $\alpha$ προς το $\beta$. Ένα κλάσμα που έχει μέσα του κλάσμα λέγεται \lexiword{σύνθετο}, «φτιαγμένο από πολλά κομμάτια».
\end{lexi}

\begin{metafrasi}
Πόσες φορές χωράει το $2$ στο $6$; & $6:2=3$ γιατί $3\cdot2=6$\\
Ομόσημοι. & πηλίκο θετικό\\
Ετερόσημοι. & πηλίκο αρνητικό\\
Διαιρώ και πολλαπλασιάζω τον αντίστροφο. & $\alpha:\beta=\alpha\cdot\dfrac{1}{\beta}$\\
Κλάσμα που έχει κλάσμα μέσα του. & $\dfrac{\frac{\alpha}{\beta}}{\frac{\gamma}{\delta}}$ (σύνθετο)\\
\end{metafrasi}

\tmima{Η θεωρία}

\begin{orismos}{πηλίκο ρητών}
Το πηλίκο της διαίρεσης $\alpha:\beta$ ή $\dfrac{\alpha}{\beta}$ ($\beta\neq0$) λέγεται και \textbf{λόγος} του $\alpha$ προς το $\beta$. Αν $\alpha:\beta=\gamma$, τότε ο $\gamma$ είναι ο μοναδικός αριθμός για τον οποίο ισχύει $\gamma\cdot\beta=\alpha$.
\end{orismos}

\begin{kanonas}{πρόσημο του πηλίκου}
Για να διαιρέσουμε δύο ρητούς, διαιρούμε τις \textbf{απόλυτες τιμές} τους και στο πηλίκο βάζουμε το πρόσημο $+$, αν είναι \textbf{ομόσημοι}, και το πρόσημο $-$, αν είναι \textbf{ετερόσημοι}.
\end{kanonas}

Για παράδειγμα $(+6):(+2)=+3$, γιατί $(+3)\cdot(+2)=+6$. Επίσης $(-8):(+4)=-2$, γιατί $(-2)\cdot(+4)=-8$, και $(-12):(-3)=+4$, γιατί $(+4)\cdot(-3)=-12$.

\begin{kanonas}{διαίρεση με τον αντίστροφο}
Για να \textbf{διαιρέσουμε} δύο ρητούς αρκεί να \textbf{πολλαπλασιάσουμε} τον διαιρετέο με τον \textbf{αντίστροφο} του διαιρέτη:
$$\alpha:\beta=\alpha\cdot\dfrac{1}{\beta}=\dfrac{\alpha}{\beta}\qquad\text{και}\qquad\dfrac{\alpha}{\beta}:\dfrac{\gamma}{\delta}=\dfrac{\alpha}{\beta}\cdot\dfrac{\delta}{\gamma}.$$
\end{kanonas}

Για παράδειγμα $\frac{2}{3}:\frac{4}{5}=\frac{2}{3}\cdot\frac{5}{4}=\frac{10}{12}=\frac{5}{6}$ και $\left(+\frac{2}{5}\right):\left(-\frac{1}{2}\right)=-\frac{2}{5}\cdot\frac{2}{1}=-\frac{4}{5}$.

\tmimasmall{Σύνθετα κλάσματα}

Ένα κλάσμα, του οποίου ένας τουλάχιστον όρος είναι κλάσμα, λέγεται \textbf{σύνθετο}. Η γραμμή του κλάσματος σημαίνει διαίρεση, άρα για να το κάνουμε \textbf{απλό} εκτελούμε τη διαίρεση:
$$\dfrac{\frac{\alpha}{\beta}}{\frac{\gamma}{\delta}}=\dfrac{\alpha}{\beta}:\dfrac{\gamma}{\delta}=\dfrac{\alpha\cdot\delta}{\beta\cdot\gamma}.$$
Αν ο παρονομαστής ή ο αριθμητής είναι ακέραιος, τον γράφουμε με παρονομαστή $1$. Αν υπάρχουν πρόσημα στους όρους, τα βγάζουμε έξω από τα κλάσματα: $\frac{-1}{2}$, $\frac{3}{-4}$ και $-\frac{3}{4}$ είναι ο ίδιος τύπος αριθμών.

\tmimasmall{Διαίρεση δεκαδικών}

Για να διαιρέσουμε δεκαδικό με δεκαδικό, πολλαπλασιάζουμε \textbf{και τον διαιρετέο και τον διαιρέτη} με την ίδια κατάλληλη δύναμη του $10$ ($10$, $100$, \dots) ώστε ο διαιρέτης να γίνει φυσικός αριθμός. Το πηλίκο δεν αλλάζει. Έτσι $4{,}2:1{,}75=420:175=2{,}4$.

\begin{kanonas}{διαίρεση με $10,\ 100,\ 1\,000,\dots$ και με $0{,}1,\ 0{,}01,\dots$}
Όταν διαιρούμε έναν δεκαδικό με $10,\ 100,\ 1\,000,\dots$ μεταφέρουμε την υποδιαστολή τόσες θέσεις \textbf{προς τα αριστερά} όσα είναι τα μηδενικά. Όταν τον διαιρούμε με $0{,}1,\ 0{,}01,\ 0{,}001,\dots$ μεταφέρουμε την υποδιαστολή τόσες θέσεις \textbf{προς τα δεξιά} όσα είναι τα δεκαδικά ψηφία του διαιρέτη.
\end{kanonas}

Για παράδειγμα $12{,}68:100=0{,}1268$ και $12{,}68:0{,}01=1\,268$.

\begin{prosoxi}
\begin{itemize}[leftmargin=1.2em,itemsep=0.3ex,topsep=2pt]
\item Δεν διαιρούμε ποτέ με το $0$. Αν $5:0=x$ τότε θα έπρεπε $x\cdot0=5$, που δεν ισχύει για κανέναν αριθμό. Το $0:5=0$ είναι σωστό.
\item Η διαίρεση δεν «μικραίνει» πάντα: $3:\frac{1}{2}=6>3$. Διαιρώντας με αριθμό μικρότερο του $1$, το πηλίκο μεγαλώνει.
\item Δεν αντιστρέφουμε και τα δύο κλάσματα: $\frac{2}{3}:\frac{4}{5}=\frac{2}{3}\cdot\frac{5}{4}$ και όχι $\frac{3}{2}\cdot\frac{5}{4}$. Αντιστρέφουμε μόνο τον \emph{διαιρέτη}.
\item Η διαίρεση δεν είναι αντιμεταθετική: $6:2=3$ ενώ $2:6=\frac{1}{3}$.
\item Στη διαίρεση με $0{,}1$ η υποδιαστολή πάει \emph{δεξιά} (το πηλίκο μεγαλώνει) και όχι αριστερά.
\end{itemize}
\end{prosoxi}

\tmima{Λυμένα παραδείγματα}

\paradeigma{Τα δοχεία του παππού Γιάννη}
\tixero $450$ λίτρα λάδι και δοχεία των $1\frac{1}{2}=\frac{3}{2}$ λίτρων.\par
\tithelo Πόσα δοχεία θα γεμίσουν;\par
\possindeo Διαιρώ: $450:\frac{3}{2}=450\cdot\frac{2}{3}=\frac{900}{3}=300$.\par
\apantisi $300$ δοχεία.\par
\elegxos $300\cdot\frac{3}{2}=450$.

\paradeigma{Πρόσημα στη διαίρεση}
\tixero Οι διαιρέσεις $(+6):(+2)$, $(-8):(+4)$, $(-12):(-3)$ και $\left(-\frac{2}{5}\right):\left(-\frac{3}{10}\right)$.\par
\tithelo Να βρω τα πηλίκα.\par
\possindeo Διαιρώ τις απόλυτες τιμές και βάζω πρόσημο: ομόσημοι $+$, ετερόσημοι $-$. Για το τελευταίο: $\frac{2}{5}\cdot\frac{10}{3}=\frac{20}{15}=\frac{4}{3}$.\par
\apantisi $+3$, $-2$, $+4$, $+\dfrac{4}{3}$.\par
\elegxos Για κάθε πηλίκο ο πολλαπλασιασμός με τον διαιρέτη δίνει τον διαιρετέο: $(-2)\cdot(+4)=-8$.

\paradeigma{Διαίρεση κλασμάτων}
\tixero Οι $\frac{2}{3}:\frac{4}{5}$ και $\left(+\frac{2}{5}\right):\left(-\frac{1}{2}\right)$.\par
\tithelo Να τις υπολογίσω.\par
\possindeo Πολλαπλασιάζω με τον αντίστροφο του διαιρέτη. Για το πρώτο ο αντίστροφος του $\frac{4}{5}$ είναι ο $\frac{5}{4}$. Για το δεύτερο ο αντίστροφος του $-\frac{1}{2}$ είναι ο $-2$.\par
\apantisi $\dfrac{2}{3}\cdot\dfrac{5}{4}=\dfrac{10}{12}=\dfrac{5}{6}$ και $\dfrac{2}{5}\cdot(-2)=-\dfrac{4}{5}$.\par
\elegxos $\frac{5}{6}\cdot\frac{4}{5}=\frac{20}{30}=\frac{2}{3}$.

\paradeigma{Σύνθετα κλάσματα}
\tixero Τα $\dfrac{\frac{3}{4}}{\frac{9}{8}}$, $\dfrac{\frac{5}{6}}{10}$ και $\dfrac{4}{\frac{2}{3}}$.\par
\tithelo Να τα γράψω ως απλά κλάσματα.\par
\possindeo Γράφω κάθε ένα ως διαίρεση: $\frac{3}{4}:\frac{9}{8}=\frac{3}{4}\cdot\frac{8}{9}=\frac{24}{36}=\frac{2}{3}$. Το δεύτερο: $\frac{5}{6}:10=\frac{5}{6}\cdot\frac{1}{10}=\frac{5}{60}=\frac{1}{12}$. Το τρίτο: $4:\frac{2}{3}=4\cdot\frac{3}{2}=6$.\par
\apantisi $\dfrac{2}{3}$, $\dfrac{1}{12}$, $6$.\par
\elegxos Στο τρίτο: $6\cdot\frac{2}{3}=4$.

\paradeigma{Παράσταση με σύνθετο κλάσμα}
\tixero Η παράσταση $\dfrac{\frac{3}{4}+\frac{1}{2}}{\frac{5}{6}-\frac{1}{3}}$.\par
\tithelo Να βρω την τιμή της.\par
\possindeo Εκτελώ τις πράξεις στον αριθμητή και στον παρονομαστή: $\frac{3}{4}+\frac{1}{2}=\frac{5}{4}$ και $\frac{5}{6}-\frac{1}{3}=\frac{5}{6}-\frac{2}{6}=\frac{3}{6}=\frac{1}{2}$. Τότε $\frac{5}{4}:\frac{1}{2}=\frac{5}{4}\cdot2=\frac{10}{4}$.\par
\apantisi $\dfrac{5}{2}$.\par
\elegxos Με δεκαδικούς: $1{,}25:0{,}5=2{,}5$.

\paradeigma{Διαίρεση δεκαδικών}
\tixero Η διαίρεση $3{,}6:0{,}45$ (κιλά αλεύρι προς κιλά ανά κέικ).\par
\tithelo Πόσα κέικ βγαίνουν;\par
\possindeo Πολλαπλασιάζω και τους δύο αριθμούς επί $100$, ώστε ο διαιρέτης να γίνει φυσικός: $360:45=8$.\par
\apantisi $8$ κέικ.\par
\elegxos $8\cdot0{,}45=3{,}6$.

\paradeigma{Διαίρεση με $10$, $100$, $0{,}1$, $0{,}01$}
\tixero Ο αριθμός $12{,}68$.\par
\tithelo Να τον διαιρέσω με $10$, $100$, $0{,}1$, $0{,}01$.\par
\possindeo Με $10$ και $100$ η υποδιαστολή πάει αριστερά κατά $1$ και $2$ θέσεις. Με $0{,}1$ και $0{,}01$ πάει δεξιά κατά $1$ και $2$ θέσεις.\par
\apantisi $1{,}268$, \ $0{,}1268$, \ $126{,}8$, \ $1\,268$.\par
\elegxos Το $12{,}68:0{,}1$ είναι «πόσα δέκατα χωράνε στο $12{,}68$», δηλαδή $126{,}8$.

\tmima{Ασκήσεις}

\zoni{Το κατάλαβα}{κανόνες και πράξεις}

\ex Κάνε τις διαιρέσεις: α) $15:(-5)$ \quad β) $(-60):(+15)$ \quad γ) $(+81):(-9)$ \quad δ) $(-490):(-7)$.

\ex Βρες τα πηλίκα: α) $\frac{2}{5}:\frac{4}{15}$ \quad β) $\left(+\frac{1}{2}\right):\left(-\frac{1}{3}\right)$ \quad γ) $4:\frac{1}{2}$ \quad δ) $\left(-\frac{4}{3}\right):(-8)$ \quad ε) $\frac{18}{5}:\frac{6}{5}$ \quad στ) $\frac{12}{21}:\frac{7}{4}$.

\ex Μετάτρεψε τα σύνθετα κλάσματα σε απλά: α) $\dfrac{\frac{7}{11}}{\frac{3}{22}}$ \quad β) $\dfrac{\frac{27}{49}}{\frac{9}{14}}$ \quad γ) $\dfrac{4}{\frac{56}{9}}$ \quad δ) $\dfrac{\frac{11}{2}}{121}$ \quad ε) $\dfrac{\frac{3}{4}}{2\frac{1}{8}}$.

\ex Υπολόγισε τα πηλίκα: α) $4{,}5:0{,}15$ \quad β) $28{,}42:49$ \quad γ) $35{,}76:1{,}5$ \quad δ) $292{,}62:29{,}262$ \quad ε) $0{,}48:0{,}6$.

\ex Υπολόγισε χωρίς κάθετη πράξη: α) $45{,}7:10$ \quad β) $3{,}2:100$ \quad γ) $0{,}5:1\,000$ \quad δ) $6{,}4:0{,}1$ \quad ε) $0{,}35:0{,}01$ \quad στ) $12:0{,}001$.

\ex Σωστό ή λάθος; α) $\alpha:\beta=\alpha\cdot\frac{1}{\beta}$. \quad β) Επιτρέπεται η διαίρεση με το $0$. \quad γ) Το πηλίκο ετεροσήμων είναι αρνητικό. \quad δ) $(-12):(-4)=-3$. \quad ε) $3:\frac{1}{3}=1$. \quad στ) Το πηλίκο δύο ρητών είναι πάντα μικρότερο από τον διαιρετέο.

\ex Υπολόγισε: α) $\dfrac{-3}{4}-\dfrac{5}{-8}+\dfrac{-1}{-2}$ \quad β) $1+\dfrac{5}{-4}-\dfrac{-1}{8}$.

\ex Βρες τον αριθμό $x$ ώστε: α) $(-3)\cdot x=12$ \quad β) $x\cdot\frac{2}{3}=4$ \quad γ) $(-5)\cdot x=-\frac{1}{2}$.

\zoni{Το χρησιμοποιώ}{μοιράζω και μετράω}

\ex \textbf{Το λάδι.} Ο παππούς Γιάννης έχει $6$ λίτρα λάδι και θέλει να το βάλει σε μπουκάλια των $\frac{3}{4}$ του λίτρου. Πόσα μπουκάλια θα γεμίσουν;

\ex \textbf{Η κορδέλα.} Η Αγγελική κόβει μια κορδέλα $4{,}5$ μ. σε κομμάτια των $0{,}15$ μ. Πόσα κομμάτια βγαίνουν;

\ex \textbf{Η τιμή.} Η γιαγιά Μυρτώ πλήρωσε $6$ € για $0{,}75$ κιλά τυρί. Πόσο κοστίζει το κιλό;

\ex \textbf{Το χρέος του ταμείου.} Το ταμείο της εκδρομής έχει υπόλοιπο $-24{,}60$ €. Τα $6$ παιδιά της παρέας μοιράζονται εξίσου το χρέος. Τι αριθμός εκφράζει το μερίδιο του καθενός;

\ex \textbf{Η ταχύτητα.} Ο Νίκος τρέχει $\frac{3}{5}$ του χιλιομέτρου σε $\frac{1}{10}$ της ώρας. Πόσα χιλιόμετρα τρέχει σε μία ώρα (αν διατηρεί τον ίδιο ρυθμό);

\zoni{Το σκέφτομαι}{γιατί το πηλίκο μεγαλώνει;}

\ex Ο Βασίλης λέει ότι η διαίρεση πάντα μικραίνει τον αριθμό. Βρες αντιπαράδειγμα και εξήγησε πότε το πηλίκο είναι μεγαλύτερο από τον διαιρετέο.

\ex Γιατί δεν επιτρέπεται η διαίρεση με το $0$; Σκέψου τι θα σήμαινε $5:0=x$ αν υπήρχε τέτοιος αριθμός.

\ex Ποιο είναι μεγαλύτερο: το $3:\frac{1}{2}$ ή το $3:\frac{1}{3}$; Εξήγησε με τη φράση «πόσες φορές χωράει».

\ex Αν $\frac{\alpha}{\beta}>0$, τι πρόσημα μπορούν να έχουν οι $\alpha$ και $\beta$; Αν $\frac{\alpha}{\beta}<0$;

\zoni{Για το σπίτι}{κόβω και μετράω}

\ex Πάρε ένα σπάγκο μήκους $1{,}2$ μ. και κόψε τον σε κομμάτια των $0{,}15$ μ. Πρόβλεψε πρώτα με διαίρεση πόσα θα βγουν και μετά έλεγξε.

\ex Γράψε τρία σύνθετα κλάσματα της επιλογής σου και μετάτρεψέ τα σε απλά.
